\documentclass[titlepage]{article}

%packages
%\usepackage{biblatex}
\usepackage[numbers]{natbib}
\usepackage{titlepic}
\usepackage{graphicx}

\title{Autonomous Car Intersection Simulation Project}
\author{
   Ivan Naumovski\\
    {inau@itu.dk}\\
    \\
   Martino Secchi\\
    {msec@itu.dk}\\
}
\titlepic{\includegraphics[scale=.3]{img/itulogo.jpg}}
\date{December 2016}


%\addbibresource{report2.bib}
\begin{document}

\begin{titlepage}
	\centering

	{\huge\bfseries Autonomous Car Intersection Simulation Project \par}
	\vspace{2cm}
	{\Large Ivan Naumovski, {inau@itu.dk}\\ Martino Secchi, {msec@itu.dk}\\\par}
	\vfill
	supervised by\par
	\large Prof.~Kasper \textsc{St\o y}

	\vfill
	\includegraphics[width=0.3\textwidth]{img/itulogo.jpg}\par\vspace{1cm}
% Bottom of the page
	{\large \today\par}
\end{titlepage}

\pagenumbering{gobble}
%\maketitle
\clearpage
\tableofcontents
\clearpage

\pagenumbering{arabic}

\abstract{
\begin{em}
This project aims to investigate some of the aspects involving traffic management with autonomous vehicles.
In particular, we want to know how it can be possible to regulate self driving cars through an intersection, and if it is possible to achieve this without a centralised controller and the use of communication.
We built a simulation environment using the Unity3D game engine. We tried different representations of intersections and designed a reactive controller.
The data we gathered described throughput, collisions and deadlocks.
The controller we built performs worse than the compared centralised solution from Texas University.
The strength is that by going pure reactive we avoid having a single point of failure planning mechanism.
\end{em}
}
\section{Introduction}
%A presentation of previous research (related work), different paradigms of control (basic theory), and our angle.

The idea of having autonomous agents has existed in science fiction literature for hundreds of years and can also be traced back to ancient mythology.
The ideas and concepts of autonomous agents are hence not something entirely new, but rather something that has been refined over the course of a lot of years.
However the technology has just recently reached a state where a lot of the underlying challenges to autonomous agents can be solved somewhat efficiently.
Self driving cars are becoming more and more a reality, and they will be an important presence in the near future.
By automising vehicle control, traffic management can arguably be made more efficient in many ways, starting with throughput but also in other aspects, like fuel management and emissions control. 

In this paper we will focus on what it takes to make autonomous automobiles efficient in passing intersections.

Initially we had a look at what other researchers had been doing the recent decade and quickly realised that there was a pattern. 
The same groups keep appearing in the literature when having the focus on autonomous vehicles and intersections.

To put everything in context we will start by presenting some of the challenges that are present when trying to efficiently solve the task of managing vehicles in an intersection.
Followed by different control paradigms. These paradigms will be presented in relation to existing solutions to match the theory to their practical implementation.
%emphasize how practical applications of these abstract paradigms can look.  


\subsection{Motivation}
The theory we have touched upon is from multiple fields of engineering. One field is within traffic engineering, namely intersection management. The other is within the field of robotics, namely robot control.
First we will briefly explain why intersection management is deemed relevant, followed by a similar section for robot control paradigms.

\subsubsection{Intersection Management}
The field of intersection management tries to deal with the challenges arising when trying to organize multiple entities to efficiently pass through a critical zone.\\
The problem to solve is to go from origin to destination while not colliding with other entities.\\

\marginnote{Determining priority}
In the real world these challenges are solved by defining a set of rules and designing the critical zones such that it is trivial to determine from where entities might enter and leave the zone and what entity has priority over the others.\\

Different templates exist for solving these tasks efficiently and these templates can differ from nation to nation.
However our starting point is rooted in the way that traffic is handled in Europe, where rules are mostly homogenous especially in terms of traffic-legislation and intersection design.\\
Traffic-legislation contains a set of rules that makes it easy to determine what entity has a higher priority within critical sections (intersections).\\
Intersections are also designed in such a way that agents have an easier task at assigning priorities. 
This includes tools such as traffic lights, road markings and/or signs.

\marginnote{Intersection types}
The intersection models that we are interested in are 4 way cross-intersections both with and without the use of traffic lights.
This is due to the fact that these types of intersections are the most common ones in our local area. Additionally, a smaller intersection still follows the same rules and can be seen as a 4 way intersection where traffic never arrives from, or travels to, one of the ways.

\subsubsection{Control Paradigms and Related Work}
We use the term agents for robots or other actors moving within an environment.\\
\marginnote{Theory}
We use the term environment as being synonymous to the world in which the agent is confined.
The following will contain short presentations of the three canonical ways of doing robot control, based on the theory presented in the textbook The Robotics Primer\citep{book:mataric}, and chapters 12 to 15 on the subject.\\

\marginnote{Reactive}
\textit{Reactive Control} (also named Sense-Act) is a paradigm that acts directly on sensory input.\\
In our research we did not stumble upon any systems using purely reactive controllers, which we found interesting. This is despite the fact that the simplicity of the interaction with the environment should make pure reactive controllers quite safe.
Simple schemes for collision avoidance are trivial to implement using pure reactive controllers.\\
An example of such a controller can be expressed quite simply using only a frontal sensor. The rule could be phrased as: "accelerate while nothing is in front of the vehicle, brake otherwise".
%while this sensor does not register anything in front of the vehicle for a certain distance accelerate else brake.\\
The strength of this paradigm is that it is highly reactive to the environment. There are no complex layers in between.\\

\marginnote{Deliberate}
In contrast to the above paradigm there exists a paradigm named \textit{deliberate control}, or alternatively Sense-Plan-Act.
Compared to the reactive control it introduces a layer of planning.
The additional step however requires more information about the environment.

This is computationally heavier since it includes sensing, planning of a step, execution of the step and starting over.
This also introduces the challenge of how to represent the environment, since a bad representation could reduce performance significantly(excessive memory usage or slow lookup times).\\

\marginnote{Hybrid}
The third canonical way of doing robot control is a hybrid of the two above, and hence it is also referred to as \textit{hybrid control}.
It resides somewhere in the middle where it has abstractions of routines using reactive controllers - but uses a higher level mechanism to combine these for advanced control schemes.\\

%We discovered multiple papers that facilitated either deliberate or hybrid approaches.

Most of the research that has been conducted on intersection management systems for autonomous vehicles tends to prefer centralised systems that can handle traffic requests.
Particularly relevant is the research of two major groups in this area: the researchers at the University of Texas at Austin, and an international group of researchers composed by members of the Massachusetts Institute of Technology (MIT), the Swiss Institute of Technology (ETHZ), and the Italian National Research Council (CNR).

At the University of Texas, Kurt Dresner and Peter Stone developed AIM, Autonomous Intersection Management, a reservation-based system built around a detailed communication protocol able to coordinate movement of self driving cars through intersections \cite{texas}.
Through a simulation they are able to demonstrate the potential of this system to outperform current intersection control mechanism: traffic lights and stop signs.
In the simulation, the intersection center is divided into a n x n grid of reservation tiles.
Through a "first come, first served" policy, approaching vehicles make a request to the system to reserve the space-time they need to cross the intersection. 
Their trajectory is then computed by the system and if the requesting vehicle at any time occupies a reservation tile that is already in use, the request is rejected.
The vehicle will then continue requesting until it can pass.

A similar approach is also followed by the international group of MIT, ETHZ and CNR researchers.
They developed a centralised slot-based intersection \cite{mit}.
In their simulations, cars adjust their velocities in approaching the intersection in order to arrive and cross at a given slot of time that is made available for them.
This system also involves communication between vehicles and a centrlised controller.\\

%In particular we were inspired by \citeauthor{texas}\citep{texas}, which explores the same area as us.\\
%Additionally, \citeauthor{texas} define certain criteria that need to be met for the interactions between intersection and multiple agents to be deemed succesful.
%They name properties such as \textit{autonomy for every agent},\textit{realistic sensor models}, \textit{deadlock avoidance},\textit{safety} and \textit{efficiency}, which we adopted into our model.

%They also mention more properties - such as protocol standardization, low communication complexity and Incremental Deployability, which we deem less important due to our approach.

\subsection{Problem Statement}
The previously obtained insights about robot control, intersection management and previous research led us to the following plan of attack:

\begin{quotation}
The goal is to explore the possibility of building a decentralized system which can guide autonomous cars through intersections. 
The current research and solutions we have found are focused around centralized systems. 

Centralized systems rely a lot on scheduling and coordination in a 'global' scale, 
while the decentralized system we envision would be built on a reactive approach which is centered around local state. 

We would like to build a prototype using the decentralized approach and compare results to an existing centralized solution.
\end{quotation}

\section{Considerations}
Post-research we were in a limbo, since the landscape of tooling that can be used for setting up experiments in this field of studies is  quite  vast.
However it was decided early in the proces that it would be mainly software simuation, due to the fact that robots have a higher cost and getting the amount of agents that makes these types of studies interesting is somewhat higher than a dozen of robots.

\subsection{Simulation Tools}
Performing simulations is not trivial. Different criteria should be met depending on the accuracy of the simulation that is to be performed.
In particular one should decide what level of abstraction fits the scenario, as 1:1 simulations are significantly harder to model, but might not necessarilly add value matching the increased complexity.\\

This introduces the question, what modelling complexity is sufficient for modelling our problem?\\

While the real world has three dimensions, we only need a two dimensional model.
The height dimension is not required since we can consider every agent to occupy a space in two dimensional space.
Restricting the simulated world to only two axis makes the model alot simpler to model.
If this was a simulation for aircrafts the height dimension would make a lot more sense to include.\\
%Another point is that we need the capability of being able to express physical shapes and collisions between these shapes.\\
%being able to apply kinematic forces to these shapes.

Two tools have been considered for doing the simulated world. One was the GazeboSimulator, which is a widely adopted tool in the robotics community.
Another was doing a simulation using game development tools, in particular Unity3D.

While both tools provide the functionality we need such as physics simulation in 2D/3D  environments, the degree of tools supported out-of-the-box varies.
Gazebo has robotic simulations as being its main function, and hence has alot of useful tools for robotics integrated into it.
Unity3D on the other hand is a general-purpose tool so it requires some tailoring to get the same functionality as Gazebo.

Unity3D however is a strong competitor due to the fact that we have prior experience using this tool. It is easy to go from concept to an actual prototype rapidly. And last but not least has a thriving community of professionals, academics and hobbyists aspiring to make the software great.\\
Our solution is built using Unity3D.\\

The choice, albeit it might seem arbitrary, is based on the fact that \textit{Gazebo} has a very steep learning curve.


\section{Implementation}
As we have settled on using Unity3D we work within its boundaries.
It is component driven, meaning every object in the world consists of a collection of components and a transform - the transform is the objects spatial information.
These components can be either colliders, renderers, or custom scripts, and objects can contain other objects - they can be composite.

The implementation was done by splitting the work in three major blocks. One was \textit{the vehicle model}, another was \textit{the intersection model} and the last \textit{the reactive controller} for the car.

\begin{itemize}
%\item Presentation of our model.
%\item Intersection w/wo traffic lights. Cars going straight vs. cars turning.
%\item Car model and sensor modelling (sick vs simplified directional).
%\item Graph for navigation.
\item Reactive controller.
\item Special rules (right hand).
\item Sensor range limitation based on braking distance considerations.
\item SimScale to RealWorldScale
\end{itemize}

\subsection{The World Model}
As observed in the real world, multiple intersection models exist, three lane and 4 lane intersections are the most common ones.
The style can be either with the roundabout approach, which essentially rules out the left turn complication or a traditional 4 lane intersection where going left crosses the opposing traffics lane.

In our model we mainly focus on a single lane 4 way intersection.\\
\textbf{insert simple intersection images - maybe graph images from unity}
%We also consider a 4 way roundabout.

We represent these intersection models using an underlying directed graph.
 A car agent in our environment is seeded with information about its origin and its destination.
These points are vertices in our graph. 
Every combination of origin and destination pair has one path.

While this does not provide alot of flexibility in straying from the path, the simplicity makes the model really easy to implement and adjust.

\subsubsection{The traffic lights}
This is an attempt to mimic traffic light regulated intersections.
The lights change at intervals chosen in regards to our experiences with intersections. 
These values can vary alot from intersection to intersection - and are highly influenced by traffic flow during the day. 
We did experiments with different interval windows.


\subsection{The Car Model}
The way the car has been modelled is using a interface which essentially contains actions such as accellerating or decellerating, braking, and turning left or right.
This abstraction simplifies the step of going from a virtual model to a physical one with robot actors.
This interface is then used by the controller. This seperation makes it simple to develop multiple types of controllers.

\subsection{The Reactive Controller}
This was developed incrementally. Initial versions were done using only collision avoidance.
The vision was to have multiple different controllers and be able to hotswap these for the simulation.




\section{Experiments}
we have some fixed parameters such as light control on or off. Mimics intersection lights.\\
Whether simulated vehicles are able to turn or not.\\
We have different max speeds - 60 or 40 km per hour approach. This translates to lower speeds during turns (roughly half).\\

tests are different combinations of above parameters. And results are measured in collisions, cars spawned, cars reaching destination, and time before deadlock.\\
Deadlock is the situation where cars have been stuck for a period of time in the `intersection zone' without entering or leaving.





\section{Discussion}
\iffalse
\begin{itemize}
\item Discuss experimental results. ~
\item Shortcomings of model.
\item Improvements for further tests.
\item traffic fairness - distribution
\end{itemize}
\fi

\subsection{Shortcomings and limitations}
 
As discussed in the previous section, the system has some clear shortcomings.
First of all, safety is not ensured since there is still occurrence of \textbf{collisions} between vehicles.
According to our observations though, most of the collisions happened in deadlock or semi-deadlock situations in which the cars were moving extremely slowly trying to pass through.
This is a condition that could be easily corrected with a more advanced sensory layer, and arguably doesn't really represent a safety hazard for potential passengers.

Another case in which collisions were relatively frequent was when cars approached the intersection at high speeds.
In 1 out of 10 tests, a collision happened in the first minute of the simulation.
Our system can handle collision avoidance in most cases, but with specific angles of impact and speed there is not enough brake distance to react in time.
\newline

Another shortcoming of our system is the abundant presence of deadlocks.
This condition is harder to solve, because it is caused by an issue underlying in the structure of the intersection.
Cars trying to pass through the same space concurrently can easily block each other, especially without a centralised controller that can sort the traffic.

\subsection{Improvements and further work}
 %sensory layer, obstacle avoidance, mimic locks with pure reactive (don't go if anyone is inside, safer, probably slower, but not in the long run)\\
%roundabouts, other paradigms that are less complex

Here we discuss about how some of the issues encountered could be approached and how further work could be structured.

To improve \textbf{safety} of the system, further improvements of the sensory layer could be performed.
The current model is rather simplistic in its implementation, especially comparing it to real life car sensor systems.
More angles and distances can be added to the sensing part, to be sure to include every possible scenario.

On the same topic, also the controller can definitely be extended to handle more advanced situations and establish more rules to follow. 
A complete navigation system that can also vary the path would be a good starting point for obtaining obstacle avoidance, but it would not be mandatory.

Obstacle avoidance could be also added as part of the current controller, allowing cars to make variations on the navigation graph currently implemented.
\newline

Another very important aspect that needs improving is the handling of \textbf{concurrency} and resulting deadlocks.
In this issue lies the strength of a centralised controller approach, that can sort out all the traffic with the focused goal of avoiding race conditions and deadlocks.

To achieve this in a decentralised approach, each agent would need to know a lot about the environment, namely the presence, direction and speed of every other agent surrounding it.
Acting on this information and through a common set of rules, each agent would then proceed with its task in a sorted way.

All of this is hardly achievable without communication.

Computer vision might be one of the technologies up to the task, along with advanced artificial intelligence techniques.
\newline

Another thing that could produce interesting and potentially better results is the use of a different and simpler model, like an intersection with a roundabout.
Roundabouts remove complexity in the handling of collision avoidance, by making the centre of the intersection not viable and thus eliminating one of the main reasons that produce deadlocks.

In a roundabout based intersection, each car has to check only two directions for collision avoidance, namely its front and its left side when entering.
Unfortunately we didn't have the time to test this approach during this project, but it would be a very interesting addition to our work.
\newline

Finally, another source of discussion is traffic distribution and fairness.
In all of our simulations, traffic was evenly generated from every direction.

This is not the case in most real intersections, where traffic usually flows in one main direction.
This could likely be a parameter in a more complete system, and performance results would probably differ from the ones registered before.
To test this behaviour would open also other angles to the problem, like for instance fairness in the handling of traffic.
When traffic is homogeneous, it is important to handle it fairly to distribute the flow evenly.
In the situation where traffic is not homogeneous however, this might not be the case, as prioritising one end or the other can potentially lead to better performance. 



\section{Conclusion}
The reactive approach can be seen as a stepping stone to achieve autonomy.
The nature of the controller makes more advanced schemes more complex to incorporate.
Our controller is an example of how far you can get with subsumption architectures - we did collision avoidance and built right hand rule onto that.

While our systems performance is comparable to AIM, it is nowhere near as good.
The throughput is lower. Our system can deadlock.
This naturally implies that we have room for improvement.
With additional work put into the controller, it should be achievable to design a controller which does not collide.

Despite the fact that it performs poorly compared to other systems, its strength is the omission of centralised control. Centralised control introduces a construct prone to single point of failure.

The vision was that pure reactive control should provide a safer way to traverse intersections with less sources for errors.


\bibliographystyle{unsrtnat}
\bibliography{report2}
%\printbibliography

\end{document}